% =====================================================================
% intro.tex — Capítulo 1: Introdução
% =====================================================================

\chapter{Introdução}
\label{cap:intro}

\section{Contexto e Motivação}
\label{sec:contexto}

A leitura e a interpretação de cartas topográficas constituem competências
fundamentais para o oficial combatente. A capacidade de visualizar, em três
dimensões, o relevo de uma região de operações é determinante para o
planejamento de manobras, o posicionamento de armas de apoio, a escolha de
itinerários de ataque e a definição de bases de fogo. Tradicionalmente, esse
treinamento é apoiado por dois recursos complementares: a carta
topográfica em papel, com curvas de nível, e o caixão de areia
físico, no qual a turma de instrução modela manualmente o terreno a partir
da carta para obter uma representação tridimensional tangível.

Embora o caixão de areia seja um instrumento didático consagrado, a
transição cognitiva entre a representação bidimensional da carta
(curvas de nível) e o relevo tridimensional construído na areia exige
considerável esforço mental e tempo de instrução. Erros de interpretação
das curvas de nível resultam em terrenos modelados que divergem da
realidade cartográfica, comprometendo a fidelidade do exercício.

Os caixões de areia com realidade aumentada — na literatura
internacional, \textit{Augmented Reality Sandboxes} — emergiram, na última
década, como uma resposta tecnológica a esse problema. Por meio da
combinação de um sensor de profundidade (tipicamente o
Microsoft Kinect) com um projetor multimídia suspenso
sobre a caixa, esses sistemas projetam, em tempo real, curvas de nível
coloridas, escalas batimétricas ou texturas de relevo diretamente sobre a
superfície da areia. À medida que o operador modifica o relevo — cavando ou
acumulando areia com as mãos — a projeção se ajusta dinamicamente,
fornecendo \emph{feedback} visual imediato e dispensando a interpretação
mental das curvas de nível.

A iniciativa pioneira do projeto SARndbox, da Universidade da
Califórnia em Davis~\cite{sarndbox}, demonstrou a viabilidade técnica da
abordagem. Adaptações posteriores expandiram o conceito para a educação
em geociências, hidrologia e cartografia. Neste trabalho, por sugestão da
orientação, evita-se o uso do anglicismo como apelido do sistema
desenvolvido, preferindo-se, ao longo de todo o texto, a denominação em
língua portuguesa ``caixão de areia com realidade aumentada''.

No contexto militar, entretanto,
três limitações dos sistemas existentes restringem sua adoção:

\begin{enumerate}
    \item A maioria dos sistemas projeta apenas curvas de nível
          decorativas ou simulações hidrológicas qualitativas, sem
          comparação com um relevo de referência específico;
    \item Não há guia visual explícito indicando onde
          cavar ou onde acumular areia para reproduzir um terreno
          alvo determinado pela cartografia oficial;
    \item As implementações de referência são frágeis quanto a
          falhas de hardware: a ausência do Kinect ou do arquivo de
          terreno provoca o encerramento abrupto do programa,
          comprometendo apresentações didáticas e demonstrações.
\end{enumerate}

Este trabalho tem origem em uma demanda operacional concreta: a
Seção de Simulação da Academia Militar das Agulhas Negras (AMAN)
mantém, em sua Área de Emissão de Ordens Tecnológicas, um caixão
de areia físico empregado rotineiramente na instrução de leitura de carta
topográfica, no planejamento de operações e em exercícios de Comando e
Estado-Maior. Nesse ambiente, o relevo é modelado manualmente pelos
instruendos a partir de cartas em papel, e a aderência do terreno
resultante à cartografia oficial é verificada apenas visualmente pelo
instrutor, sem qualquer retorno quantitativo durante a modelagem. A
Seção de Simulação identificou nesse processo uma oportunidade de
modernização: dotar o caixão de areia já existente de um mecanismo de
\emph{feedback} automático, capaz de comparar o relevo modelado com um
Modelo Digital de Elevação de referência sem substituir a mesa física
nem alterar a rotina de instrução já consolidada. Esse é o contexto
institucional que orienta os requisitos e as decisões de projeto
apresentadas neste documento.

\section{Definição do Problema}
\label{sec:problema}

Diante das limitações expostas, este Projeto Final de Curso propõe a
seguinte questão de pesquisa:

\begin{quote}
\itshape
``Como projetar e implementar um sistema de Realidade Aumentada para
caixão de areia que (a) compare quantitativamente, em tempo real, o relevo
modelado com um Modelo Digital de Elevação (MDE) cartográfico de
referência, (b) forneça \emph{feedback} visual inequívoco indicando onde
adicionar ou remover areia para minimizar essa diferença, e (c) opere de
forma resiliente, sem qualquer falha catastrófica, mesmo na ausência do
sensor de profundidade ou do arquivo de terreno?''
\end{quote}

\section{Objetivos}
\label{sec:objetivos}

\subsection{Objetivo geral}
\label{subsec:obj-geral}

Desenvolver um sistema funcional de Caixão de Areia com Realidade
Aumentada que projete sobre uma caixa física de areia uma
malha discretizada de quadrados coloridos comparando, célula a
célula, o relevo medido por um sensor de profundidade com o relevo de
referência extraído de um arquivo GeoTIFF, e que opere em modalidade
\emph{plug \& play} com fallback automático para um emulador
interativo.

\subsection{Objetivos específicos}
\label{subsec:obj-especificos}

\begin{enumerate}
    \item Projetar uma arquitetura de software monolítica
          modularizada de três camadas (Hardware, Lógica, Dados),
          orquestrada por uma máquina de estados explícita;
    \item Implementar o módulo \texttt{KinectSensor} com mecanismo de
          \textit{cascading fallback} (PyKinect2~$\to$~Open3D~$\to$~freenect~$\to$~Simulação);
    \item Implementar o motor matemático \texttt{motor\_caixao\_areia.py}
          contemplando ajuste de plano por SVD, base ortonormal por
          Gram-Schmidt, transformação afim 4$\times$4 e projeção
          \textit{pinhole} de Tsai;
    \item Implementar o adaptador \texttt{AdaptadorMDE} para leitura de
          GeoTIFF via \texttt{rasterio} com interpolação bilinear, e
          fallback para superfície gaussiana sintética;
    \item Implementar discretização espacial em malha
          regular, com agregação de pontos por célula (filtragem
          estatística do ruído do sensor) e renderização por
          \texttt{cv2.fillPoly};
    \item Implementar um emulador interativo controlado por
          \textit{mouse} que permita cavar e preencher a areia virtual
          quando o Kinect não estiver disponível;
    \item Validar todo o motor matemático através de uma suíte de
          testes unitários automatizados, seguindo os
          princípios do \textit{Test-Driven Development} (TDD);
    \item Documentar o projeto em três níveis: documentação oficial
          acadêmica, README de uso prático e relatório de PFC
          em conformidade com as normas da ABNT.
\end{enumerate}

\section{Contribuição e Justificativa}
\label{sec:contribuicao}

A contribuição central deste trabalho reside na integração inédita de
três elementos: (i) o rigor matemático de uma calibração
geométrica completa (SVD + Gram-Schmidt + Tsai), (ii) a regra
de negócio quantitativa de comparação com tolerância numérica explícita
contra um MDE oficial, e (iii) uma arquitetura resiliente que
prioriza a confiabilidade operacional na presença de falhas de hardware
ou de dados. Diferentemente de implementações que utilizam apenas mapas
de cores qualitativos, o sistema proposto produz uma decisão
operacional para cada célula da mesa: cavar, preencher ou manter.

A justificativa militar é direta: a doutrina nacional prescreve o uso
do caixão de areia para o ensino e o planejamento tático em todas as
escolas de formação. Reduzir o tempo necessário à modelagem do relevo
e aumentar a fidelidade cartográfica do exercício produz ganho direto
em horas de instrução e em qualidade do treinamento.

\section{Organização do Documento}
\label{sec:organizacao}

Os capítulos seguintes acompanham o próprio \textit{pipeline} do
sistema: da teoria (Capítulo~\ref{cap:fundamentacao}), passando pelo
posicionamento deste trabalho frente a implementações correlatas
(Capítulo~\ref{cap:relacionados}), e da arquitetura
de \textit{software} (Capítulo~\ref{cap:arquitetura}) à formalização
matemática da calibração e da projeção (Capítulo~\ref{cap:motor}), à
renderização e ao emulador interativo (Capítulo~\ref{cap:rendering}) e,
por fim, aos resultados experimentais e à conclusão. O sumário e a
lista de figuras, nas páginas pré-textuais, detalham a numeração
completa de seções; o Apêndice~A reúne o roteiro de demonstração
utilizado na defesa do PFC.
